\clearpage
\renewcommand{\sectioncolor}{alert}
\section{What is actually still missing}
\markboth{What is still missing}{}

The v1.0 list had five books. Four were already yours. \textbf{Two remain.}

\begin{center}\small
\begin{tabular}{@{}c>{\raggedright\arraybackslash}p{3.6cm}>{\raggedright\arraybackslash}p{5.2cm}>{\raggedright\arraybackslash}p{6.3cm}@{}}
\toprule
\rowcolor{alert!12}
\textbf{\#} & \textbf{\color{alert}Gap} & \textbf{\color{alert}Get} & \textbf{\color{alert}Why}\\
\midrule
\rowcolor{bio!8}
\textbf{1} & \textbf{Markov-chain mixing times} &
 \textbf{Levin \& Peres}, \textit{Markov Chains and Mixing Times} \newline{\scriptsize (the authors host a free PDF)} &
 The only thing that explains §6 \emph{quantitatively}. You now have Le Gall and Bressloff for
 stochastic processes, but neither treats mixing time --- which is the exact quantity that makes
 sampling $Z$ fail.\\
\textbf{2} & \textbf{Free-energy methods in simulation} &
 \textbf{Frenkel \& Smit}, \textit{Understanding Molecular Simulation} &
 \textbf{You have GROMACS 2023.3 compiled and no text for the mathematics it implements.} Umbrella
 sampling, thermodynamic integration, free-energy perturbation --- the practical route from $Z$ to a
 binding affinity.\\
\midrule
\multicolumn{4}{@{}l}{\textbf{\color{slate}Optional, only if you pursue the Tier 3 problem}}\\
--- & HEOM / strong-coupling open systems & Breuer \& Petruccione, \textit{The Theory of Open Quantum
 Systems} & Goes beyond N\&C ch.~8's Kraus formalism into master equations proper.\\
\bottomrule
\end{tabular}
\end{center}

\begin{haveit}[title={Books v1.0 told you to buy that you already own}]
\begin{center}\small
\begin{tabular}{@{}>{\raggedright\arraybackslash}p{7.3cm}>{\raggedright\arraybackslash}p{9.1cm}@{}}
\toprule
\rowcolor{bio!12}
\textbf{v1.0 said ``buy''} & \textbf{Where it already was}\\
\midrule
Chandler or Kardar, statistical mechanics & \textbf{Feynman, \textit{Statistical Mechanics}} --- \bk{Module\_8}. Arguably better than either for your purposes, because ch.~2--3 go straight to the density matrix\\
Grimmett \& Stirzaker, probability & \textbf{Le Gall} (409 pp) and \textbf{Lauritzen} --- \bk{Module\_8}\\
Reed \& Simon, functional analysis & \textbf{Axler \textit{MIRA}} (430 pp) --- \bk{Module\_5}; and \textbf{Bru \& Pedra} for the operator-algebra end --- \bk{QBIT-HIS}\\
Hall or Sakurai, quantum mechanics & \textbf{Sakurai \& Napolitano 3ed} (+ 2 solution manuals), Shankar, Griffiths, Townsend, Binney \& Skinner --- \bk{TensoQi101}\\
\bottomrule
\end{tabular}
\end{center}
\end{haveit}

\clearpage
\section{A twelve-week plan, revised}
\markboth{Twelve-week plan}{}

\begin{center}\small
\begin{tabular}{@{}>{\raggedright\arraybackslash}p{1.4cm}>{\raggedright\arraybackslash}p{5.5cm}>{\raggedright\arraybackslash}p{4.3cm}>{\raggedright\arraybackslash}p{5.1cm}@{}}
\toprule
\rowcolor{bio!14}
\textbf{\color{bio}Weeks} & \textbf{\color{bio}Read} & \textbf{\color{bio}Do} & \textbf{\color{bio}Question to answer in writing}\\
\midrule
\textbf{1--2} & Axler \textit{LADR} §5--7, §10.A; then \textbf{N\&C §2.1.8} &
 Prove $\operatorname{Tr}(T^N)=\lambda_+^N+\lambda_-^N$ &
 Why does $e^{-\beta\hat H}$ need the spectral theorem to exist?\\
\textbf{3} & --- & \textbf{Run the notebook}; exercises 1--3 &
 At what $N$ does brute force stop on your laptop? Extrapolate to $N=60$.\\
\rowcolor{tenso!7}
\textbf{4} & \textbf{Shankar §21.2} &
 Re-derive the Ising $Z$ his way &
 What exactly is being Wick-rotated, and why does dimension go up by one?\\
\textbf{5--6} & \textbf{Feynman} ch.~1 (\textit{The Partition Function}), ch.~2, ch.~5 &
 Do §5.2 by hand; read §5.4 (Onsager) without trying to follow every step &
 Why is 1D solvable and 2D a landmark?\\
\rowcolor{qbit!7}
\textbf{7--8} & \textbf{PBoC2 ch.~6 \textit{Entropy Rules!}} &
 Work the gene-expression model &
 Which of Part IX Fig.~2's six quantities can you now compute for a real cell?\\
\textbf{9} & \textbf{N\&C §2.4, §11.3, ch.~8} &
 Write $\rho=e^{-\beta\hat H}/Z$ and verify $\mathrm{Tr}\,\rho=1$ &
 Where does positivity get used, exactly?\\
\textbf{10} & \textbf{Axler \textit{MIRA}} ch.~1--3 (or Le Gall) &
 Verify $\mu(dx)=e^{-\beta U}dx/Z$ is a probability measure &
 In what sense does a positive operator \emph{determine} a measure?\\
\rowcolor{bio!8}
\textbf{11--12} & \textbf{Levin \& Peres} ch.~1--4 \newline{\scriptsize (the one book to acquire)} &
 Notebook exercise 4: autocorrelation time vs $\xi$ &
 \textbf{Restate §6 in the language of mixing times.}\\
\bottomrule
\end{tabular}
\end{center}
\vspace{-2pt}
\begin{center}\begin{minipage}{0.93\textwidth}\footnotesize\color{slate}
\textbf{Table 5.}\; \textbf{One book to acquire; eleven weeks from your own shelves.} At the end you
can read Part IX §7 on Boltzmann generators and know precisely what problem it solves.
\end{minipage}\end{center}

\begin{plain}
The plan is one week longer than it needs to be on purpose. Weeks 4 and 5--6 cover the same
calculation twice, from quantum mechanics and from statistical mechanics. That redundancy is the
point: it is the difference between knowing a derivation and knowing an object.
\end{plain}
